\section{Research priorities and conjectures not yet resolved}
\label{sec:priorities}

\subsection{First priority: extract and formalize the analytic selector}

The strongest completed transfer is Theorem~\ref{thm:analytic-torsion-selector}.
It relies on a local root-counting argument and has a tractable
formalization path:
\begin{enumerate}
\item Strengthen the one-coordinate update lemma to retain
\(\sum_i(a_i-c_i)_+\), then telescope across coordinates.
\item Formalize uniform total-degree tails on a fixed torsion grid and
algebraic approximation of finitely many coefficients and witnesses.
\item Apply the finite-intersection argument on the finite grid.
\item Package determinant, maximal-minor, and normalized-length corollaries
as a reusable specialization API.
\end{enumerate}
No new global RH machinery is needed. The current 103-module import
closure is an engineering starting point; reducing incidental imports
would improve portability.

The result's explicit order bound is poor for computation. For
\(p=2,s_0=1,\rho=1/2,B=1,\epsilon=1/4\), the displayed sufficient
choice is
\[
 H=\lceil3e^{24}\rceil,\qquad \text{maximum torsion order }2^H.
\]
This is roughly \(H=7.95\cdot10^{10}\), before exponentiating again.
It proves uniform existence, not a feasible enumeration algorithm.
Improving the harmonic-pigeonhole bound for particular coefficient
classes or determinant structures is a concrete secondary problem.

\subsection{Second priority: attack the \texorpdfstring{\(A_5\)}{A5} arithmetic estimate}

The exact proposed nonabelian target is
\[
 R_W(x)\le C_\epsilon x^{-1/2+\epsilon},
 \qquad W=(21,1,0,1,1).
\]
The positive detector includes all prime powers and handles inertia.
The analytic conclusion, the pole normalization, the real residue signs,
the chosen field, and the limited character optimization have been proved.
They should be treated as finished infrastructure.

The unresolved task is to construct arithmetic objects whose index or
determinant encodes this \(W\), with matching main terms and errors
small enough to preserve the square-root saving. A reasonable first
attempt would make the degree-5, degree-6, and degree-12 permutation
realizations interact with a common lattice comparison. It would need
to prove integrality, sufficient rank, a compatible long-form estimate,
and local twist control. This is a proposed research route, not an
existing construction.

The mean-square lower bound \(7\) in
Theorem~\ref{thm:a5-blindness} gives a useful early failure test:
any proposal that uses only finitely many fixed congruence classes to
represent \(W\) exactly cannot work. A growing-modulus proposal must
supply new uniform estimates and explain how it retains the source's
scale. One cannot infer failure of every RH-based argument from this
limited obstruction.

The author announcement itself identifies extended RH as a possible
further direction~\cite{alpoge-remark}. It provides motivation, not
the missing nonabelian construction. Approximate Artin holomorphy
and zero-free results for families are also substantial existing work,
but they do not provide the fixed-field, full-critical-line statement
required by this target~\cite{ltz2024}.

\subsection{Third priority: finish a finite class-group classification}

Theorem~\ref{thm:genus-four} strengthens a specific gap in the prior
analysis. It is the closest direction to a complete further
classification theorem because a published finite global range and an
existing census are already available. A next computation should apply
the actual inequality
\[
 x\le4R(H)\,M_{2^{t-\dim H}}(\log x)^6
\]
for useful subgroups \(H\), rather than only the four rounded thresholds.
Optimizing ramified representatives can reduce \(R(H)\).
The remaining discriminants would then need a complete, independently
checkable enumeration or a further theoretical exclusion.

The present report does not count that remaining set, estimate its
computational cost, or verify all 432 published entries. Any claim that
the classification has now been completed would go beyond both the
proofs and the computations supplied here.

\subsection{Why BSD, Goldfeld, Vandiver, and general Iwasawa claims do not follow}

For BSD or Goldfeld-type nonvanishing, a selector must be connected to
the correct analytic or algebraic object. A nonzero initial witness,
integral interpolation or a presentation, compatibility of local
conditions, and any required density statement remain independent
inputs. The scalar-inertia obstruction shows that even the source's
packaged Selmer hypothesis fails for the unmodified elliptic Tate module.
A selector cannot remove a forced sign zero or create a nonzero
Euler-system class from an identically zero input.

For Vandiver-type or higher \(K\)-group vanishing statements, controlling
a denominator at varying ramified twists does not show that the desired
fixed component vanishes. Normalized lengths can decrease simply
because the coefficient-field ramification increases. One would need
a fixed-component comparison and a uniform positive lower gap for
nonzero lengths, or a different argument proving exact vanishing.
The previous reports already identify related barriers~\cite{prior}.

Thus these are candidate domains for the determinant machinery, not
conjectures resolved here. A useful follow-up should name one exact
representation, module, and specialization problem and verify its
hypotheses before claiming a new arithmetic consequence.

\section{Verification record and use of the archive}
\label{sec:repro}

\subsection{Claim and evidence ledger}

\begin{center}
\small
\begin{tabularx}{\textwidth}{@{}p{49mm}YY@{}}
\toprule
Claim & Evidence supplied & Limit\\
\midrule
Scope of upstream theorem &
Pinned \file{Main.lean}, \file{Statements.lean}, source hashes,
reported independent checking pages. &
No new local Lean build or kernel replay.\\
\addlinespace
Analytic torsion theorem &
Complete proof including positive excess, approximation, uniform tails,
and the finite-intersection step. &
Finite Python examples do not certify the infinite theorem.\\
\addlinespace
\(A_5\) transfer criterion &
Complete transform/residue proof and exact finite-group and polynomial
certificate. &
The prime-sum inequality remains a hypothesis.\\
\addlinespace
Genus-field exclusion theorem &
Complete class-group/norm proof, original prime-bound constants, exact
rational endpoint checks. &
Published census is an external input; no full survivor enumeration.\\
\addlinespace
Novelty relative to packets &
Inspection of all five prior archives, pinned manifest, detailed
comparison notes. &
No exhaustive literature-priority determination.\\
\bottomrule
\end{tabularx}
\end{center}

The exact finite checks have distinct purposes. The \(A_5\) program
constructs the group and coset actions, verifies the character identities,
local cycle factors, variance and the Galois certificate. The genus
program checks rational inequalities, table ranges, the demonstration
primes and discriminant. The torsion program checks illustrative finite
valuation computations and a harmonic bound. All decisive arithmetic in
the first two programs is exact. Approximate decimals in the JSON are
display aids, not the basis of comparisons.

\subsection{Contents and commands}

The archive contains \file{report.tex}, \file{report.pdf}, the included
section files, Python certificates, their JSON/CSV results,
source-audit and prior-report manifests, review notes, and a small
selection of upstream source files with their license notices.
The upstream excerpt is evidence, not a standalone buildable copy of
\repo.

Run \texttt{bash code/reproduce.sh} from the archive root to regenerate
the finite certificates and rebuild the PDF when the dependencies are
available. Python 3 is sufficient for the \(A_5\) and genus scripts;
the torsion illustrations also require SymPy. A standard LaTeX installation
with \file{latexmk} and \file{pdflatex} builds the report. The source audit
accepts the path to a separately cloned repository at the pinned revision.
The included SHA-256 manifest checks archive members; it does not attest
to their mathematical correctness.

